\documentclass[11pt]{article}

% ---------- Packages ----------
\usepackage[margin=1in]{geometry}
\usepackage[T1]{fontenc}
\usepackage{lmodern}
\usepackage{microtype}
\usepackage{amsmath, amssymb, amsthm}
\usepackage{graphicx}
\usepackage{booktabs}
\usepackage{siunitx}
\usepackage{caption}
\usepackage{float}
\usepackage{xcolor}
\usepackage{tcolorbox}
\usepackage[colorlinks=true, linkcolor=blue!50!black, urlcolor=blue!50!black]{hyperref}

\captionsetup{font=small, labelfont=bf, margin=1cm}
\sisetup{group-digits=false, detect-weight=true, mode=text}
\setlength{\parskip}{0.5em}
\setlength{\parindent}{0pt}

\theoremstyle{plain}
\newtheorem{theorem}{Theorem}

% Placeholder box for notebook scans -- swap for \includegraphics once ready
\newcommand{\scanplaceholder}[1]{%
  \begin{tcolorbox}[colback=gray!5, colframe=gray!50, arc=2pt, boxrule=0.6pt,
                    halign=center, valign=center, height=4cm]
    \textit{\color{gray!70!black}#1}
  \end{tcolorbox}}

% ---------- Title ----------
\title{\textbf{The Ballot Theorem}\\[0.3em]
       \large An Experimental Verification via Monte Carlo Simulation}
\author{Samuel Shindich}
\date{}

\begin{document}
\maketitle

% =====================================================================
\section{Introduction}

The Ballot Theorem, first published by W. A. Whitworth in 1878 and later rediscovered by
Joseph Louis Bertrand in 1887, is as follows:

\begin{theorem}[Ballot Theorem]
Suppose candidate $A$ receives $a$ votes and candidate $B$ receives $b$ votes,
with $a > b$, and all $\binom{a+b}{a}$ orderings of the ballots are equally
likely. Then the probability that $A$ is strictly ahead of $B$ throughout the
entire count is
\[
  P = \frac{a-b}{a+b}.
\]
\end{theorem}

In this project, the Ballot Theorem is verified experimentally using Monte Carlo
simulations on a variety of $(a,b)$ pairs.

% =====================================================================
\section{Proof (Using the Reflection Principle)}

For the $i$-th ballot counted, let $X_i = +1$ if it is an $a$-vote and
$X_i = -1$ if it is a $b$-vote, and let
\[
  S_0 = 0, \qquad S_k = X_1 + \cdots + X_k \quad (1 \le k \le a+b),
\]
so $S_k$ is $A$'s lead after $k$ ballots. We model $(S_k)$ as a simple symmetric
random walk, i.e.\ the $X_i$ are independent with
$P(X_i = 1) = P(X_i = -1) = \tfrac12$. Every path with $S_{a+b} = a-b$ has
exactly $a$ up-steps and $b$ down-steps, so each such path has the same
probability $1/2^{a+b}$. Hence, conditioned on $S_{a+b} = a-b$, all
$\binom{a+b}{a}$ orderings of the ballots are equally likely, which is exactly
the setting of the Ballot Theorem.

Let $N_n(x,y)$ denote the number of paths of length $n$ from height $x$ to
height $y$. Let $N^{0}_n(x,y)$ denote the number of such paths that intersect
the horizontal axis, and let $N^{+}_n(x,y)$ denote the number of such paths that
stay strictly above $0$.

By the definition of conditional probability,
\[
  P\bigl(S_k > 0 \ \forall k \in \{1,\dots,a+b\} \bigm| S_{a+b} = a-b\bigr)
  = \frac{P\bigl(S_k > 0 \ \forall k \in \{1,\dots,a+b\},\ S_{a+b} = a-b\bigr)}
         {P\bigl(S_{a+b} = a-b\bigr)}.
\]
The denominator is $\frac{1}{2^{a+b}} N_{a+b}(0,\,a-b)$. For the numerator, a
successful path must begin with an $a$-vote, taking it to height $1$; from there,
any path counted by $N^{+}_{a+b-1}(1,\,a-b)$ completes it. So the numerator is
$\frac{1}{2^{a+b}} N^{+}_{a+b-1}(1,\,a-b)$, and the factors of $2^{-(a+b)}$
cancel:

\begin{align*}
P\bigl(S_k > 0 \ \forall k \bigm| S_{a+b} = a-b\bigr)
  &= \frac{N^{+}_{a+b-1}(1,\, a-b)}{N_{a+b}(0,\, a-b)}
   = \frac{N_{a+b-1}(1,\, a-b) - N^{0}_{a+b-1}(1,\, a-b)}{\binom{a+b}{a}} \\[0.5em]
  &= \frac{N_{a+b-1}(1,\, a-b) - N_{a+b-1}(-1,\, a-b)}{\binom{a+b}{a}}
   = \frac{\binom{a+b-1}{a-1} - \binom{a+b-1}{a}}{\binom{a+b}{a}} \\[0.5em]
  &= \frac{\dfrac{(a+b-1)!}{(a-1)!\,b!} - \dfrac{(a+b-1)!}{a!\,(b-1)!}}
          {\dfrac{(a+b)!}{a!\,b!}}
   = \frac{\dfrac{(a-b)\,(a+b-1)!}{a!\,b!}}{\dfrac{(a+b)!}{a!\,b!}}
   = \frac{a-b}{a+b}. \qquad \blacksquare
\end{align*}

% =====================================================================
\section{Experiment}

The 28 $(a,b)$ pairs selected for this experiment are listed in
Figure~\ref{fig:pairs} below.

\begin{figure}[H]
  \centering
  \begin{tabular}{c l}
    \toprule
    $a$ & $b$ \\
    \midrule
    3  & 1 \\
    5  & 3 \\
    10 & 9,\; 5 \\
    15 & 14,\; 10,\; 5 \\
    20 & 19,\; 15,\; 10 \\
    25 & 24,\; 20,\; 15 \\
    30 & 29,\; 25,\; 20 \\
    35 & 34,\; 30,\; 25 \\
    40 & 39,\; 35,\; 30 \\
    45 & 44,\; 40,\; 35 \\
    50 & 49,\; 45,\; 40 \\
    \bottomrule
  \end{tabular}
  \caption{The 28 $(a,b)$ pairs tested, grouped by $a$.}
  \label{fig:pairs}
\end{figure}

For any given pair $(a,b)$ such that $a \geq b$, I created a list containing
$a$ ones and $b$ zeros. Then, for a single trial, the order of the list's
elements was randomized, and the counts of both $A$-votes and $B$-votes were set
to $0$. The list was then iterated through from the beginning, adding a vote to
the $A$-count every time a $1$ was encountered, and to the $B$-count when a $0$
was encountered. If at any point the number of $B$-votes was equal to or
surpassed the number of $A$-votes (after the first vote had been cast), then the
trial failed. Otherwise, if the list was fully iterated through without the
$B$-votes catching up to the $A$-votes, then the trial succeeded.

For all $(a,b)$ pairs, the above trial was run $N = 10{,}000$ times, tallying
the number of successful trials. The resulting fraction is listed in the
\emph{experimental} column (representing the experimental probability),
alongside the pair's theoretical probability given by the Ballot Theorem,
$\frac{a-b}{a+b}$. The difference between the two probabilities is also listed
to highlight how far off the results were after $10{,}000$ trials.

% =====================================================================
\section{Results}

The full table of results is shown in Figure~\ref{fig:results} below. Note that
across all pairs, the largest difference (in absolute value) between the
experimental and theoretical results was $0.0146$ for $(a,b) = (15,5)$, and this
was the only pair with a difference greater than $0.01$. Thus, the experimental
results are very close to the theoretical ones given by the Ballot Theorem,
indicating that the Ballot Theorem holds in practice.

\begin{figure}[H]
  \centering
  \small
  \begin{tabular}{
      S[table-format=2.0]
      S[table-format=2.0]
      S[table-format=1.5]
      S[table-format=1.4]
      S[table-format=-1.5]
      S[table-format=-1.5]}
    \toprule
    {$a$} & {$b$} & {Theoretical} & {Experimental} & {Difference} & {Relative diff.} \\
    \midrule
     3 &  1 & 0.5     & 0.4993 & -0.0007  & -0.0014  \\
     5 &  3 & 0.25    & 0.2532 &  0.0032  &  0.0128  \\
    \addlinespace
    10 &  9 & 0.05263 & 0.0523 & -0.00033 & -0.00627 \\
    10 &  5 & 0.33333 & 0.3331 & -0.00023 & -0.00069 \\
    \addlinespace
    15 & 14 & 0.03448 & 0.0344 & -0.00008 & -0.00232 \\
    15 & 10 & 0.2     & 0.2032 &  0.0032  &  0.016   \\
    \bfseries 15 & \bfseries 5 & \bfseries 0.5 & \bfseries 0.5146
               & \bfseries 0.0146 & \bfseries 0.0292 \\
    \addlinespace
    20 & 19 & 0.02564 & 0.0257 &  0.00006 &  0.00234 \\
    20 & 15 & 0.14286 & 0.146  &  0.00314 &  0.02198 \\
    20 & 10 & 0.33333 & 0.3334 &  0.00007 &  0.00021 \\
    \addlinespace
    25 & 24 & 0.02041 & 0.0211 &  0.00069 &  0.03381 \\
    25 & 20 & 0.11111 & 0.114  &  0.00289 &  0.02601 \\
    25 & 15 & 0.25    & 0.2513 &  0.0013  &  0.0052  \\
    \addlinespace
    30 & 29 & 0.01695 & 0.0179 &  0.00095 &  0.05605 \\
    30 & 25 & 0.09091 & 0.0901 & -0.00081 & -0.00891 \\
    30 & 20 & 0.2     & 0.205  &  0.005   &  0.025   \\
    \addlinespace
    35 & 34 & 0.01449 & 0.0154 &  0.00091 &  0.0628  \\
    35 & 30 & 0.07692 & 0.0777 &  0.00078 &  0.01014 \\
    35 & 25 & 0.16667 & 0.1659 & -0.00077 & -0.00462 \\
    \addlinespace
    40 & 39 & 0.01266 & 0.0131 &  0.00044 &  0.03476 \\
    40 & 35 & 0.06667 & 0.0684 &  0.00173 &  0.02595 \\
    40 & 30 & 0.14286 & 0.1444 &  0.00154 &  0.01078 \\
    \addlinespace
    45 & 44 & 0.01124 & 0.0114 &  0.00016 &  0.01423 \\
    45 & 40 & 0.05882 & 0.058  & -0.00082 & -0.01394 \\
    45 & 35 & 0.125   & 0.1225 & -0.0025  & -0.02    \\
    \addlinespace
    50 & 49 & 0.0101  & 0.0096 & -0.0005  & -0.0495  \\
    50 & 45 & 0.05263 & 0.0521 & -0.00053 & -0.01007 \\
    50 & 40 & 0.11111 & 0.1101 & -0.00101 & -0.00909 \\
    \bottomrule
  \end{tabular}
  \caption{Theoretical vs.\ experimental probabilities that $A$ stays strictly
           ahead, with $N = 10{,}000$ trials per pair. The largest absolute
           difference, at $(15,5)$, is shown in bold.}
  \label{fig:results}
\end{figure}

% =====================================================================
\section{Applications}

Any Catalan number $C_n$ can be derived from the special case
$(a,b) = (n+1,\, n)$ using a rearrangement of the Ballot Theorem in terms of $N^{+}_{2n+1}(0,1)$:
\[
N^{+}_{2n+1}(0,1)
  = P(a \text{ ahead of } b) \cdot N_{2n+1}(0,1)
  = \frac{(n+1)-n}{(n+1)+n} \cdot \binom{2n+1}{n}
  = \frac{1}{2n+1}\binom{2n+1}{n}
  = C_n.
\]

Of course, the Catalan numbers can also be derived directly using the reflection
principle, but it is interesting that the Ballot Theorem yields a general
formula for them.

\end{document}
